% SEGMENT OLP-0353-S01
% Part: lambda-calculus
% Chapter: introduction
% Section: primitive-recursion

\documentclass[../../../include/open-logic-section]{subfiles}

\begin{document}

\olfileid{lam}{int}{pr}
\olsection{\usetoken{S}{lambda definable} சார்புகள் முதன்மை மீள்வரையறையின் கீழ் மூடியவை}

முதன்மை மீள்வரையறையைப் பொறுத்தவரை, இறுதியாக நாம் ஓரளவு வேலை செய்ய
வேண்டும். படிப்படியாகச் செல்ல வேண்டியிருக்கும். முன்புபோல, சார்புகள் $g$
மற்றும்~$h$ ஐ முறையே !!{lambda define} செய்யும் சொற்கள் $G$ மற்றும்~$H$
ஏற்கெனவே உள்ளன என்ற கருதுகோளில், பின்வருமாறு வரையறுக்கப்படும் சார்பு~$f$ ஐ
!!{lambda define}s ஒரு சொல்~$F$ நமக்கு வேண்டும்.
% SOURCE CORRECTION TA-LAM-007: the recursive-step argument is x, not the parameter vector z.
\begin{align*}
f(0, \vec z) & = g(\vec z) \\
f(x+1, \vec z) & = h(x, f(x,\vec z), \vec z).
\end{align*}
% SOURCE CORRECTION TA-LAM-008: F is the sought term; H already names the step-function term.
எனவே பொதுவாக, லாம்டா சொற்கள் $G'$ மற்றும்~$H'$ கொடுக்கப்பட்டால், ஒவ்வோர்
இயல் எண்~$n$ க்கும் பின்வருவன அமையுமாறு ஒரு சொல்~$F$ ஐக் கண்டால் போதும்:
\begin{align*}
F(\num{0}, \vec z) & \equiv G(\vec z) \\
F(\overline{n+1}, \vec z) & \equiv H(\num{n}, F(\num{n}, \vec z), \vec z)
\end{align*}
$G'$ மற்றும்~$H'$ ஆகியவை $g$ மற்றும்~$h$ ஐ !!{lambda define} செய்கின்றன
என்பதால், $\vec z$ க்குப் பதிலாக எண்ணுருக்கள் $\num{\vec m}$ ஐ
இடைநிறுத்தும்போதெல்லாம், $F(\num{n+1}, \num{\vec m})$ சரியான விடைக்கு
இயல்பாக்கப்படும்.

ஆனால் இதற்கு, பின்வருவனவற்றை நிறைவுசெய்யும் சொல்~$F$ ஒன்றைக் கண்டாலே போதும்:
% SOURCE CORRECTION TA-LAM-009: v represents the preceding value as a function of the parameter vector, so it is applied only to that vector.
\begin{align*}
F(\num 0) & \equiv G \\
F(\overline {n+1}) & \equiv H(\num{n},F(\num{n}))
\intertext{ஒவ்வோர் இயல் எண் $n$ க்கும்; இங்கு}
G & = \lambd[\vec z][G'(\vec z)] \text{ மற்றும்}\\
H(u,v) & = \lambd[\vec z][H'(u,v(\vec z),\vec z)].
\end{align*}
வேறு சொற்களில், லாம்டா கலனத்தின் உத்திகளைக் கொண்டு கூடுதல் அளவுருக்கள்
$\vec z$ பற்றிக் கவலைப்படுவதைத் தவிர்க்கலாம்---அவை லாம்டா குறிமானத்திற்குள்
உட்கவரப்படுகின்றன.

சொல் $F$ ஐ வரையறுப்பதற்கு முன்பு, வரிசைப்படுத்தப்பட்ட ஜோடிகளைக் கையாளும்
ஒரு செயல்முறை நமக்குத் தேவை. அதை அடுத்த துணைத்தேற்றம் தருகிறது.

\begin{lem}
லாம்டா சொற்களின் ஒவ்வொரு ஜோடி $M$, $N$ க்கும்
$D(M,N)(\num{0}) \red M$ மற்றும் $D(M,N)(\num{1}) \red N$ ஆக அமையும்
ஒரு லாம்டா சொல் $D$ உள்ளது.
\end{lem}

\begin{proof}
முதலில், லாம்டா சொல் $K$ ஐப் பின்வருமாறு வரையறுக்க:
\[
K(y) = \lambd[x][y].
\]
வேறு சொற்களில், $K$ என்பது $\lambd[y][\lambd[x][y]]$ என்ற சொல். இதை
வேறுவிதமாகப் பார்த்தால், ஒவ்வொரு $M$ க்கும் $K(M)$ என்பது எந்த உள்ளீட்டிலும்
~$M$ ஐத் திருப்பித்தரும் மாறிலிச் சார்பு.

இப்போது $D(x,y,z)$ ஐ $D(x,y,z) = z (K(y))x$ என வரையறுக்க. அப்போது
\begin{align*}
D(M,N,\num 0) & \red \num 0 (K(N)) M \red M \text{ மற்றும்}\\
D(M,N,\num 1) & \red \num 1 (K(N)) M \red K(N) M \red N,
\end{align*}
தேவையானவாறே கிடைக்கின்றன.
\end{proof}

$D(M,N)$ என்பது ஜோடி $\tuple{M, N}$ ஐப் பிரதிநிதித்துவப்படுத்துகிறது என்பதே
கருத்து. $P$ அத்தகைய ஒரு ஜோடியைப் பிரதிநிதித்துவப்படுத்துவதாகக் கொண்டால்,
$P(\num 0)$ மற்றும் $P(\num 1)$ ஆகியவை இடது, வலது வீழ்த்தல்கள் $(P)_0$,
$(P)_1$ ஐப் பிரதிநிதித்துவப்படுத்துகின்றன. பின்னுள்ள இக்குறிமானங்களையே
பயன்படுத்துவோம்.

\begin{lem}
!!{lambda definable} சார்புகள் முதன்மை மீள்வரையறையின் கீழ் மூடியவை.
\end{lem}

\begin{proof}
ஏதேனும் சொற்கள் $G$, $H$ கொடுக்கப்பட்டால், ஒவ்வோர் இயல் எண்~$n$ க்கும்
\begin{align*}
F(\num{0}) & \equiv G \\
F(\num{n+1}) & \equiv H(\num{n}, F(\num{n}))
\end{align*}
ஆக அமையும் ஒரு சொல் $F$ ஐக் கண்டறிய முடியும் என்று காட்ட வேண்டும். தோராயமான
கருத்து, எண்ணுருக்களை மீள்செயல்படுத்திகளாகப் பயன்படுத்தி
\[
\tuple{\num 0, F(\num 0)}, \tuple{\num 1, F(\num 1)}, \dots,
\]
என்ற \emph{ஜோடிகளின்} தொடர்களைக் கணிப்பதாகும். முதல் ஜோடி
$\tuple{\num 0, G}$ மட்டுமே என்பதைக் கவனிக்க. ஒரு ஜோடி
$\tuple{\num{n}, F(\num{n})}$ கொடுக்கப்பட்டால், அடுத்த ஜோடி
$\tuple{\num{n+1}, F(\num{n+1})}$ ஆனது
$\tuple{\num{n+1}, H(\num{n}, F(\num{n}))}$ குச் சமானமாக இருக்க வேண்டும்.
இந்த ஒருபடிநிலை மாற்றத்தைச் செய்யும் லாம்டா சொல்~$T$ ஒன்றை வடிவமைப்போம்.

விவரங்கள் பின்வருமாறு. $T(u)$ ஐ
\[
T(u) = \tuple{S((u)_0), H((u)_0,(u)_1)}.
\]
என வரையறுக்க. இப்போது எந்த எண்~$n$ க்கும்
\[
T(\tuple{\num{n}, M}) \red \tuple{\num{n+1}, H(\num{n}, M)}.
\]
என்பதை எளிதாகச் சரிபார்க்கலாம். மேலே கூறியபடி, $G$, $H$ கொடுக்கப்பட்டால்,
~$F(u)$ ஐ
\[
F(u) = (u(T,\tuple{\num 0, G}))_1.
\]
என வரையறுக்க. வேறு சொற்களில், உள்ளீடு~$\num{n}$ மீது $F$ ஆனது
$\tuple{\num 0, G}$ இல் $T$ ஐ $n$ முறை மீளச்செயல்படுத்தி, பின்னர் இரண்டாம்
கூற்றைத் திருப்பித்தருகிறது. தொடக்கத்தில் நமக்கு
\begin{enumerate}
\item $\num{0} (T, \tuple{\num 0, G}) \equiv \tuple{\num{0}, G}$
\item $F(\num{0}) \equiv G$
\end{enumerate}
கிடைக்கின்றன. $n$ மீதான தொகுத்தறிதலால், ஒவ்வோர் இயல் எண்ணுக்கும் பின்வருவன
அமையும் என்று காட்டலாம்:
\begin{enumerate}
\item $\num{n+1}(T, \tuple{\num{0}, G}) \equiv \tuple{\num{n+1}, F(\num{n+1})}$
\item $F(\num{n+1}) \equiv H(\num{n}, F(\num{n}))$
\end{enumerate}
இரண்டாம் கூறுக்கு,
\begin{align*}
F(\num{n+1}) & \red (\num{n+1}(T, \tuple{\num 0, G}))_1 \\
& \equiv (T(\num{n} (T, \tuple{\num 0, G})))_1 \\
& \equiv (T(\tuple{\num{n}, F(\num{n})}))_1 \\
& \equiv (\tuple{\num{n+1}, H(\num{n}, F(\num{n}))})_1 \\
& \equiv H(\num{n}, F(\num{n})).
\end{align*}
இங்கு கடைசிக்கு முந்திய வரியில் தொகுத்தறிதல் கருதுகோளைப் பயன்படுத்தியுள்ளோம்.
முதல் கூறுக்கு,
\begin{align*}
\num{n+1} (T, \tuple{\num 0, G}) &
\equiv T(\num{n} (T, \tuple{\num 0, G})) \\
& \equiv T( \tuple{\num{n}, F(\num{n})}) \\
& \equiv \tuple{\num{n+1}, H(\num{n}, F(\num{n}))} \\
& \equiv \tuple{\num{n+1}, F(\num{n+1})}.
\end{align*}
இங்கு கடைசி வரியில் இரண்டாம் கூறைப் பயன்படுத்தியுள்ளோம். ஆகவே
$F(\num 0) \equiv G$ என்றும், ஒவ்வொரு $n$ க்கும்
$F(\num {n+1}) \equiv H(\num n, F(\num{n}))$ என்றும் காட்டியுள்ளோம்;
இதுவே நமக்குத் தேவையானது.
\end{proof}

\end{document}
